\section{Introduction}\label{rl-0001}

I don't want to write an introduction right now, but here's some stuff to mention.
\begin{itemize}
  \item We use the Chow groups of Kresch, but \(\hat A_*^T\mathcal Y\) does not mean the Chow groups of Edidin-Graham-Totaro that appear in his Definition 2.1.4 (iii). We need this group in order to make sense of the Euler classes of vector bundle stacks.
  \item The perfect obstruction theory of \(M^T\) is relative to the map \(M^T\to \mathcal Y^T\), \textit{not} the map \(M^T\to \mathcal Y\). We get this obstruction theory by restricting \(E\) to \(M^T\) and taking the fixed portion. This is the primary source of difficulty. We must find a way to relate the virtual fundamental classes \([M]^{vir}\) and \([M^T]^{vir}\) by studying the map \(\eta:\mathcal Y^T\to \mathcal Y\) between the bases, which is impossible via traditional localization theory since the main localization theorem fails for Artin stacks. We fix this by defining a residue map from classes of \(\mathcal Y\) to classes of \(\mathcal Y^T\). There is no canonical obstruction theory for \(\eta\) so the virtual pullback \(\eta^!\) (if it exists) depends on the choice of POT chosen, but this dependency can be fixed by capping with an appropriate inverse Euler class.
\end{itemize}

Each section begins with its own introduction explaining its context within the broader goal of the paper and is followed by subsections containing the actual math.

\bigskip

Throughout \(\mathcal X\) is an algebraic stack of finite type over a field \(k\) with an action of \(T\) a torus. We'll always denote by \(\eta:\mathcal X^T\to \mathcal X\) the canonical map from the Romagny fixed locus to \(\mathcal X\), and \(\mathcal Z\subseteq \mathcal X^T\) will be an open and closed sublocus. The restriction of \(\eta\) to \(\mathcal Z\) will also just be denoted \(\eta\). We'll need the following hypotheses.

\noindent
\begin{itemize}
  \item \(T\) is a split torus. We need this so that \(e_T(V)\) is a unit whenever \(V\) is a finite-rank moving \(T\)-equivariant locally free sheaf on \(\mathcal Z\subseteq \mathcal X^T\).
  \item \(\mathcal X\) is stratified by global quotients. Need this for Kresch's Chow groups and also for Manolache's virtual pullback.
  \item \(\mathcal Z\) has the resolution property. We need this for the existence of a POT for \(\eta:\mathcal Z\to \mathcal X\). Note this implies \(\mathcal Z = [U/\GL_n]\) with \(U\) some quasi-affine by \cite{totaro_ResolutionPropertySchemes2004}.
  \item \(\mathcal X^T\) is algebraic. This is not automatic, and is equivalent to \(\eta:\mathcal X^T\to \mathcal X\) being representable by algebraic spaces.
\end{itemize}

This means \(A^T_*(\mathcal Z) = A_*(\mathcal Z\times BT)\) exists with no further hypotheses.
